\documentclass[12pt,a4paper]{report}
\usepackage{fontspec}
\setmainfont{Latin Modern Roman}
\usepackage[vietnamese]{babel}
\usepackage[a4paper,top=2.3cm,bottom=2.3cm,left=3cm,right=2cm]{geometry}
\usepackage{amsmath,amssymb,graphicx,booktabs,float,caption,hyperref,listings,xcolor,tabularx}
\usepackage{microtype}
\setlength{\parskip}{0.4em}
\captionsetup{font=small,labelfont=bf}
\graphicspath{{../results/figures/}}
\lstset{basicstyle=\ttfamily\small,breaklines=true,frame=single,columns=fullflexible,
  backgroundcolor=\color{gray!8},language=Python,showstringspaces=false}
\hypersetup{hidelinks,pdftitle={Backpropagation và Khởi tạo tham số trong Deep Learning}}
\addto\captionsvietnamese{\renewcommand{\contentsname}{MỤC LỤC}}
\addto\captionsvietnamese{\renewcommand{\listfigurename}{DANH MỤC HÌNH}}
\addto\captionsvietnamese{\renewcommand{\listtablename}{DANH MỤC BẢNG}}
\addto\captionsvietnamese{\renewcommand{\bibname}{TÀI LIỆU THAM KHẢO}}
\addto\captionsvietnamese{\renewcommand{\appendixname}{PHỤ LỤC}}

\begin{document}
\pagenumbering{roman}

\chapter*{LỜI CẢM ƠN}
Nhóm chúng em xin cảm ơn các thầy cô đã hướng dẫn và cung cấp kiến thức nền tảng về Deep Learning trong suốt học phần. Báo cáo được thực hiện trên tinh thần trung thực học thuật: mọi số liệu đều đến từ các lần chạy thực tế, mọi kết luận đều giới hạn trong phạm vi thí nghiệm đã thực hiện.

\chapter*{DANH MỤC TỪ VIẾT TẮT}
\begin{tabularx}{\textwidth}{lX}
\toprule
Viết tắt & Giải thích \\
\midrule
MLP & Multi-Layer Perceptron (mạng nơ-ron nhiều lớp kết nối đầy đủ) \\
ReLU & Rectified Linear Unit \\
SGD & Stochastic Gradient Descent \\
RMS & Root Mean Square \\
VJP & Vector--Jacobian Product \\
Var & Variance (phương sai) \\
NaN & Not a Number \\
\bottomrule
\end{tabularx}

\tableofcontents
\listoftables
\listoffigures

\chapter*{LỜI MỞ ĐẦU}
\addcontentsline{toc}{chapter}{LỜI MỞ ĐẦU}
\pagenumbering{arabic}

\section*{1. Bối cảnh và tính cấp thiết}
Huấn luyện mạng nơ-ron là bài toán tối ưu hoá tham số $\theta$ theo hàm mất mát $L(\theta)$. Hai câu hỏi nền tảng quyết định việc mạng có học được hay không: gradient $\nabla_\theta L$ được tính thế nào khi mạng có hàng triệu tham số, và tham số nên được khởi tạo ở giá trị nào. Hai câu hỏi này liên quan trực tiếp đến hiện tượng vanishing và exploding gradient, vốn là rào cản lớn khi huấn luyện mạng sâu.

\section*{2. Vấn đề nghiên cứu}
Backpropagation là ứng dụng có hệ thống của chain rule, nhưng nó chỉ tính gradient. Gradient đi qua nhiều lớp là tích của nhiều Jacobian, và scale ban đầu của tham số quyết định tín hiệu có co lại, phình ra hay được giữ ổn định hay không. Vấn đề nghiên cứu: \emph{khởi tạo tham số ảnh hưởng thế nào đến lan truyền gradient và khả năng học của mạng sâu?}

\section*{3. Câu hỏi nghiên cứu}
\begin{itemize}
  \item \textbf{RQ1.} Backpropagation tính gradient như thế nào, và vì sao gradient ở lớp đầu có thể vanish hoặc explode?
  \item \textbf{RQ2.} Cách khởi tạo weight (Zero, Random nhỏ, Random lớn, Xavier, He) gây ra các thất bại nào trong cùng một kiến trúc?
  \item \textbf{RQ3.} Xavier và He có vượt trội nhau một cách có ý nghĩa thống kê trong thiết lập này không, và gradient ổn định có đảm bảo generalization không?
\end{itemize}

\section*{4. Mục tiêu và nhiệm vụ}
Mục tiêu: trình bày backpropagation và khởi tạo tham số một cách có chứng cứ, bằng cả giải thích toán học và thực nghiệm nhỏ trên Fashion-MNIST. Nhiệm vụ: (i) xây dựng backpropagation thủ công và kiểm tra bằng autograd và sai phân hữu hạn; (ii) đo gradient và activation variance tại khởi tạo cho năm cách khởi tạo; (iii) đánh giá huấn luyện, độ sâu và số seed.

\section*{5. Đối tượng và phạm vi}
Đối tượng: gradient, activation và quỹ đạo huấn luyện của mạng MLP với activation ReLU. Phạm vi: MLP $784\to128\times10\to10$, Fashion-MNIST subset 5.000/1.000/2.000, SGD cố định $\eta=0{,}05$, 15 epoch, seed 42 (và 5 seed cho Xavier, He). Không khảo sát CNN, Transformer, Adam hay BatchNorm.

\section*{6. Phương pháp nghiên cứu}
\begin{itemize}
  \item \textbf{Nghiên cứu tài liệu} về backpropagation, chain rule và khởi tạo tham số.
  \item \textbf{Thực nghiệm có kiểm soát}: chỉ đổi cách khởi tạo weight, giữ mọi yếu tố khác cố định.
  \item \textbf{Đo lường}: gradient RMS, activation variance tại khởi tạo; train/validation/test loss và accuracy theo epoch.
\end{itemize}

\section*{7. Bố cục báo cáo}
Chương 1 trình bày cơ sở lý thuyết và nghiên cứu liên quan. Chương 2 mô tả thiết kế thực nghiệm. Chương 3 trình bày triển khai và kết quả. Phần Kết luận trả lời các câu hỏi nghiên cứu, nêu hạn chế và hướng phát triển. Phụ lục chứa mã nguồn trích dẫn.

\chapter{CƠ SỞ LÝ THUYẾT VÀ NGHIÊN CỨU LIÊN QUAN}

\section{Bài toán huấn luyện mạng nơ-ron}
Với mạng $f_\theta$, dữ liệu $(x,y)$ và hàm mất mát $L$, huấn luyện là tìm $\theta^\ast = \arg\min_\theta \mathbb E[L(f_\theta(x),y)]$. Gradient descent cập nhật $\theta_{\text{new}}=\theta_{\text{old}}-\eta\nabla_\theta L$. Gradient là vector các đạo hàm riêng và chỉ hướng tăng nhanh nhất của $L$; dấu trừ đưa tham số đi ngược hướng đó.

\section{Chain rule và Jacobian}
Với đồ thị tính toán $z=g(u)$, $L=f(z)$: $\partial L/\partial u = (\partial L/\partial z)(\partial z/\partial u)$. Khi biến là vector, mỗi phép biến đổi có một Jacobian. Gradient qua nhiều lớp là tích của các Jacobian này.

\section{Backpropagation}
Forward pass tính đầu ra và lưu các giá trị trung gian. Backward pass nhận $\delta^{(l)}=\partial L/\partial h^{(l)}$ và tính $\delta^{(l-1)} = J_l^\top\delta^{(l)}$, $\partial L/\partial W^{(l)} = \delta^{(l)}(h^{(l-1)})^\top$. Backprop không lập ma trận Jacobian đầy đủ mà chỉ thực hiện tích vector--Jacobian (VJP); autograd của PyTorch làm việc này cho từng nút. Backprop chỉ tính gradient; optimizer mới cập nhật tham số.

\section{Vanishing và exploding gradient}
Nếu mỗi lớp co norm tín hiệu theo hệ số $0{,}5$, sau 20 lớp hệ số là $0{,}5^{20}\approx9{,}5\times10^{-7}$ (vanishing). Nếu khuếch đại theo hệ số $1{,}5$, sau 20 lớp hệ số khoảng $3\times10^3$ (exploding). Với sigmoid và tanh, bão hoà làm đạo hàm về 0. Với ReLU, miền âm làm second moment $\mathbb E[a^2]$ giảm xấp xỉ một nửa mỗi lớp.

\section{Khởi tạo tham số}
\begin{itemize}
  \item \textbf{Zero:} các neuron cùng lớp giống hệt nhau; đây là thất bại đối xứng (\emph{failure to break symmetry}). Vấn đề nằm ở weight; bias khởi tạo 0 vẫn chấp nhận được.
  \item \textbf{Random} $\mathcal N(0,\sigma^2)$: $\sigma$ nhỏ dẫn tới vanishing; $\sigma$ lớn dẫn tới exploding.
  \item \textbf{Xavier/Glorot} (Normal): $\operatorname{Var}(W)=2/(n_{\text{in}}+n_{\text{out}})$, giả định activation tuyến tính quanh 0 hoặc đối xứng quanh 0 (như tanh). Mục tiêu là giúp duy trì variance, không phải đảm bảo tuyệt đối.
  \item \textbf{He/Kaiming}: $\operatorname{Var}(W)=2/n_{\text{in}}$, bù việc ReLU làm giảm $\mathbb E[a^2]$. Đại lượng được điều chỉnh là $\mathbb E[a^2]$, không phải variance thuần vì ReLU không zero-centered.
\end{itemize}

\section{Nghiên cứu liên quan và khoảng trống}
Các nghiên cứu nền tảng gồm Rumelhart và cộng sự về backpropagation~\cite{rumelhart}, LeCun và cộng sự về các mẹo huấn luyện~\cite{lecun}, Glorot và Bengio về Xavier~\cite{glorot}, và He và cộng sự về He initialization cho ReLU~\cite{he}. Phần lớn công trình này chứng minh lý thuyết hoặc đo trên benchmark lớn. Báo cáo này đóng góp một thí nghiệm kiểm soát chặt, cô lập đúng biến khởi tạo, trên mạng 10 lớp, với đo gradient và activation variance tại khởi tạo.

\chapter{THIẾT KẾ THỰC NGHIỆM}

\section{Dữ liệu}
Fashion-MNIST~\cite{fmnist} với subset 5.000 train, 1.000 validation và 2.000 test. Ảnh 28$\times$28 được làm phẳng thành 784 chiều và chuẩn hoá theo mean/std tính trên tập train.

\section{Kiến trúc và huấn luyện}
\begin{table}[H]
\centering
\begin{tabular}{ll}
\toprule
Thành phần & Giá trị \\
\midrule
Kiến trúc & MLP $784\to128\times10\to10$, ReLU \\
Optimizer & SGD thuần, $\eta=0{,}05$, batch 128, 15 epoch, không momentum, không scheduler \\
Loss & CrossEntropy \\
Bias & Khởi tạo 0 trong mọi cấu hình \\
Seed & 42; 5 seed $\{42,\dots,46\}$ cho Xavier và He \\
Biến duy nhất & Cách khởi tạo weight \\
\bottomrule
\end{tabular}
\caption{Thiết lập thực nghiệm cố định.}
\label{tab:setup}
\end{table}

\section{Biến và chỉ số}
Gradient RMS: $\sqrt{\mathrm{mean}(g^2)}$ trên một batch đầu tiên, đo trước mọi cập nhật. Activation variance đo tương tự. Train/validation/test loss và accuracy đo theo epoch; test chỉ dùng đánh giá cuối.

\section{Giao thức}
Năm cấu hình khởi tạo chạy cùng seed. Xavier và He chạy thêm 5 seed để ước lượng biến thiên. Thí nghiệm độ sâu so sánh Random, Xavier và He ở nhiều độ sâu với số epoch ít hơn, nên số tuyệt đối chỉ dùng để so sánh tương đối.

\chapter{TRIỂN KHAI VÀ KẾT QUẢ}

\section{Triển khai}
Backpropagation được cài đặt thủ công bằng NumPy và được kiểm tra bằng PyTorch autograd và sai phân hữu hạn; sai số tối đa ở mức $10^{-9}$. Phần huấn luyện dùng PyTorch với seed cố định.

\section{Kết quả tại khởi tạo}
\begin{table}[H]
\centering
\begin{tabular}{lccc}
\toprule
Khởi tạo & Gradient RMS lớp 1 & Test accuracy & Hành vi \\
\midrule
Zero & $0$ & 10,0\% & Thất bại đối xứng \\
Random nhỏ ($\sigma=0{,}01$) & $1{,}1\times10^{-13}$ & 10,0\% & Vanishing \\
Random lớn ($\sigma=1{,}0$) & $1{,}9\times10^{7}$ & 10,0\% (loss NaN) & Exploding $\to$ NaN \\
Xavier (Normal) & $4{,}4\times10^{-4}$ & 68,2\% & Không vanish/explode rõ rệt \\
He (Kaiming) & $1{,}1\times10^{-2}$ & 67,4\% & Không vanish/explode rõ rệt \\
\bottomrule
\end{tabular}
\caption{Kết quả tại khởi tạo và test accuracy sau 15 epoch (seed 42).}
\label{tab:init}
\end{table}

\begin{figure}[H]
\centering
\includegraphics[width=0.85\textwidth]{gradient_heatmap_slide.png}
\caption{Gradient RMS ($\log_{10}$) theo lớp và cách khởi tạo, đo tại khởi tạo trên một batch.}
\label{fig:heat}
\end{figure}

\begin{figure}[H]
\centering
\includegraphics[width=0.8\textwidth]{activation_variance_slide.png}
\caption{Activation variance tại khởi tạo theo lớp (forward pass), một batch, trước mọi cập nhật.}
\label{fig:var}
\end{figure}

\section{Quỹ đạo huấn luyện}
\begin{figure}[H]
\centering
\includegraphics[width=0.8\textwidth]{deep_demo_accuracy_slide.png}
\caption{Validation accuracy theo epoch cho năm cách khởi tạo cùng kiến trúc.}
\label{fig:acc}
\end{figure}
Zero và Random nhỏ dừng ở 10\%. Random lớn không hội tụ. Xavier và He đạt khoảng 67--68\% test accuracy. Dao động validation accuracy đáng kể vì dùng vanilla SGD với $\eta=0{,}05$, không momentum và không scheduler, trên mạng 10 tầng.

\section{Xavier so với He}
Trên 5 seed, test accuracy trung bình $\pm$ độ lệch chuẩn là $69{,}1\%\pm4{,}7$pp (Xavier) và $70{,}5\%\pm5{,}7$pp (He). Khoảng mean $\pm$ 1 std chồng lấn đáng kể. Gradient RMS lớp 1 của He lớn hơn Xavier khoảng 35 lần, nhất quán trên mọi seed, nhưng gradient lớn hơn không đồng nghĩa với kết quả tốt hơn. Với số seed hiện tại, chưa đủ cơ sở xếp hạng Xavier và He.

\section{Độ sâu}
\begin{figure}[H]
\centering
\includegraphics[width=0.75\textwidth]{depth_comparison_v2_slide.png}
\caption{Gradient RMS lớp 1 tại khởi tạo theo độ sâu, cho Random, Xavier và He.}
\label{fig:depth}
\end{figure}
Khi độ sâu tăng, gradient của Random và Xavier suy giảm; He giữ scale tốt hơn trong thiết lập này. Ổn định gradient không đồng nghĩa với generalization tốt: ở độ sâu 12, He có train loss thấp nhưng test loss cao. Kết luận chỉ đúng với MLP, ReLU, Fashion-MNIST và SGD đã thử.

\section{Thảo luận}
Ba thất bại ban đầu được phân biệt rõ: Zero-init thất bại vì đối xứng, Random nhỏ vì vanishing, Random lớn vì exploding; NaN chỉ là triệu chứng của mất ổn định số học. Trong mạng rất sâu (trên khoảng 50 lớp), initialization thường cần kết hợp với Residual connections và Normalization (LayerNorm, BatchNorm) để kiểm soát gradient flow; báo cáo này không thực nghiệm các kỹ thuật đó.

\chapter*{KẾT LUẬN}
\addcontentsline{toc}{chapter}{KẾT LUẬN}

\section*{1. Kết quả đạt được}
\begin{enumerate}
  \item Backpropagation được trình bày như chain rule áp dụng ngược qua đồ thị tính toán; gradient descent là quy tắc cập nhật dùng gradient đó.
  \item Trong thiết lập đã thử, zero-init gây thất bại đối xứng, random nhỏ gây vanishing, random lớn gây exploding.
  \item Xavier và He tránh được các thất bại trên trong thiết lập này.
\end{enumerate}

\section*{2. Trả lời câu hỏi nghiên cứu}
\textbf{RQ1} được trả lời bằng phần lý thuyết và kiểm tra backprop. \textbf{RQ2} được trả lời bằng bảng kết quả tại khởi tạo. \textbf{RQ3}: chưa đủ cơ sở xếp hạng Xavier và He; gradient ổn định không đảm bảo generalization tốt.

\section*{3. Hạn chế}
\begin{itemize}
  \item Dataset là subset nhỏ, chỉ một kiến trúc MLP với ReLU.
  \item Optimizer cố định (SGD, $\eta=0{,}05$), không khảo sát tương tác với learning rate hay optimizer khác.
  \item Gradient và activation variance đo tại khởi tạo trên một batch.
  \item Xavier và He có 5 seed; Zero, Random nhỏ và Random lớn chỉ có 1 seed.
  \item Thí nghiệm độ sâu dùng số epoch ít hơn; số tuyệt đối có thể khác.
\end{itemize}

\section*{4. Hướng phát triển}
Chạy nhiều seed hơn để đánh giá khác biệt Xavier và He; thử thêm Adam và BatchNorm; kết hợp initialization với Residual connections trên mạng rất sâu.

\chapter*{TÀI LIỆU THAM KHẢO}
\addcontentsline{toc}{chapter}{TÀI LIỆU THAM KHẢO}
\begin{thebibliography}{9}
\bibitem{rumelhart} D. E. Rumelhart, G. E. Hinton, R. J. Williams. Learning representations by back-propagating errors. \emph{Nature}, 323(6088):533--536, 1986.
\bibitem{lecun} Y. LeCun, L. Bottou, G. B. Orr, K.-R. Müller. Efficient BackProp. In \emph{Neural Networks: Tricks of the Trade}, 1998/2012.
\bibitem{glorot} X. Glorot, Y. Bengio. Understanding the difficulty of training deep feedforward neural networks. \emph{AISTATS}, 2010.
\bibitem{he} K. He, X. Zhang, S. Ren, J. Sun. Delving Deep into Rectifiers. \emph{ICCV}, 2015. arXiv:1502.01852.
\bibitem{goodfellow} I. Goodfellow, Y. Bengio, A. Courville. \emph{Deep Learning}. MIT Press, 2016.
\bibitem{fmnist} H. Xiao, K. Rasul, R. Vollgraf. Fashion-MNIST. arXiv:1708.07747, 2017.
\bibitem{paszke} A. Paszke et al. PyTorch: An Imperative Style, High-Performance Deep Learning Library. \emph{NeurIPS}, 2019.
\end{thebibliography}

\appendix
\chapter*{PHỤ LỤC A. MÃ NGUỒN TRÍCH DẪN}
\addcontentsline{toc}{chapter}{PHỤ LỤC A. Mã nguồn trích dẫn}
\begin{lstlisting}[caption={Tính gradient RMS tại khởi tạo (trích, experiments/run\_deep\_demo.py)}]
loss = criterion(model(x), y)
loss.backward()
grad_rms = [p.grad.pow(2).mean().sqrt().item()
            for n, p in model.named_parameters() if "weight" in n]
\end{lstlisting}

\chapter*{PHỤ LỤC B. LỆNH TÁI LẬP KẾT QUẢ}
\addcontentsline{toc}{chapter}{PHỤ LỤC B. Lệnh tái lập kết quả}
\begin{lstlisting}[language=bash]
python experiments/run_manual_backprop.py
python experiments/run_gradient_check.py
python experiments/run_deep_demo.py
python experiments/run_multiseed.py
python experiments/run_depth_experiment_v2.py
\end{lstlisting}

\end{document}
